\chapter{Conclusion}
\label{chap:conclusion}

This thesis asked whether post-COVID ECB monetary easing transmitted more strongly through housing-credit and financial-intermediation channels than through compensation-linked real-economy variables. Within the monthly reduced-form local-projection framework, the answer is a qualified yes. Easing-signed ECB policy news moves house-purchase lending growth clearly at short and medium horizons, and lending margins move with it. Whether this is stronger than transmission to compensation depends on the measure: the euro-area wage tracker responds weakly and briefly, while the German industry wage bill responds durably. The estimates pool easing and tightening surprises across 2005 to 2025, so they do not isolate the post-COVID cycle, which motivates the question rather than defining the sample. The finding concerns where policy news shows up most clearly in aggregate data. It does not show that easing benefits financial markets more than the real economy, and it is narrower than a broad asset-market result.

\section{Summary of Findings}

The full interpretation is developed in \Cref{chap:discussion}. This section restates the main responses in compact form. House-purchase lending growth is the clearest response at short and medium horizons: it reacts on impact, peaks at h6, accumulates through h12, and gives back about a quarter of that accumulation by h24 while remaining positive, with an exclusion of zero that every interval method confirms through h12, which makes it the core empirical contribution. Lending spreads corroborate this reading, as mortgage and NFC spreads are positive and cumulatively durable while broad credit quantities and ECB assets are weaker. The real DAX response is negative at medium and long horizons. It is read as information-effect sensitivity rather than a simple asset-price channel, and therefore bounds broad asset-price claims without displacing the housing-credit evidence. The euro-area compensation proxies respond less consistently. The ECB Wage Tracker excluding one-offs moves at short and intermediate horizons but loses h24 cumulative exclusion. The German industry wage-bill proxy responds positively and durably, though it is a national and sectoral series. The comparison with compensation therefore depends on the measure: housing credit responds more clearly than the euro-area wage tracker, but not more clearly than the German industry wage bill.

\section{Interpretation Boundary}

These results should be read as reduced-form evidence of heterogeneous monetary transmission. They support a clear short- and medium-horizon response of house-purchase lending growth and a sustained lending-margin response alongside it. They do not support household welfare conclusions, inequality decomposition, redistribution claims, a definitive verdict on household income incidence, or the claim that monetary easing benefits financial markets more than the real economy. The responses are measured in different units and on different samples, so they show where policy news is visible, not who gains from it. This boundary is a source of the design's credibility. The model uses easing-signed high-frequency ECB surprises and dynamic impulse responses to show where policy news stays most visible, and it asks the estimates to identify no more than the design permits.
